
        \documentclass[fleqn]{article}      
        \usepackage{amsmath}
        \usepackage{fontspec}
        \usepackage{kotex} % 한국어 지원  

        \begin{document}      
        쉬움: \\\\  
1. 적분이란 무엇인가요? \\\\  
2. 다음 중 적분의 기호는 무엇인가요? (A) ∑ (B) ∫ (C) Δ (D) √ \\\\  
3. 적분을 통해 구할 수 있는 것은 무엇인가요? (A) 면적 (B) 부피 (C) 둘 다 (D) 아무것도 아님 \\\\  
4. 다음 중 적분을 사용하지 않는 경우는? (A) 곡선 아래 면적 구하기 (B) 함수의 최대값 찾기 (C) 물체의 이동 거리 구하기 (D) 함수의 평균값 구하기 \\\\  
5. ∫x dx의 결과는 무엇인가요? (A) x² + C (B) x + C (C) x³ + C (D) 1/x + C \\\\  

중간: \\\\  
1. ∫(3x²) dx의 값을 구하세요. \\\\  
2. 적분의 정의를 설명하세요. \\\\  
3. ∫(2x + 3) dx의 결과를 구하세요. \\\\  
4. 적분이 중요한 이유는 무엇인가요? \\\\  
5. ∫(x³) dx의 결과는 무엇인가요? \\\\  

어려움: \\\\  
1. ∫(sin(x)) dx의 결과를 구하세요. \\\\  
2. ∫(e^x) dx의 결과는 무엇인가요? \\\\  
3. ∫(1/x) dx의 결과를 설명하세요. \\\\  
4. ∫(x² + 2x + 1) dx의 결과를 구하세요. \\\\  
5. ∫(x^4) dx의 결과는 무엇인가요? \\\\  

      
        \end{document} 
    